% =====================================================================
%  Chapter 29 opener -- "Chapter at a Glance", three reading paths + the
%  ladder staircase, THE running example (the 9-vertex graph G: two triangles
%  B1 = 1-2-3, B2 = 4-5-6 joined by a bridge 3-4, with a bridge pendant path
%  4-7-8-9), the bridge tree B1-B2-7-8-9, the block-cut tree, the two traps
%  (a bridge is not an articulation point; a block is not a 2-edge-component),
%  the who-owns-what comparison against Ch.12 (diameter), Ch.13 (tree DP),
%  Ch.17 (DSU), Ch.28 (isomorphism, the previous chapter), Ch.26 (virtual
%  trees) and Ch.31 (MST, the chapter after next), the recognition guide, a
%  constraint -> technique table, and notation.
%  Exports fig:c31-graph, fig:c31-bridgetree, fig:c31-taxonomy.
%  \term the NEW objects "bridge", "2-edge-connected component", "bridge tree",
%  "articulation point", "block", "block-cut tree", "functional graph", "rho",
%  "cactus" only -- diameter/tin/DSU are indexed by earlier chapters.
% =====================================================================

\begin{chaptertheorybox}{Chapter at a Glance}
Every chapter up to this one was handed a \emph{tree} and asked a question about
it.  Chapter~28 (\chref{c22}) was the last: it compared shapes, but only once you
already knew the input was acyclic and connected.  This chapter inverts the
promise.  You are handed a \emph{general} graph -- it may have cycles, cut
vertices, bridges, even out-degree-one arrows -- and the first job is to
\emph{extract the tree that is hiding inside it}, then run the tree machinery you
already own on that extracted object.

A \term{bridge} is an edge whose removal disconnects the graph; delete all of
them and the survivors split into \term{2-edge-connected components}, which glue
back together along the bridges into a tree -- the \term{bridge tree}.  A
\term{articulation point} is a vertex whose removal disconnects the graph; cut
along all of them and the graph splits into \term{blocks} (maximal
2-vertex-connected pieces, plus bridges and isolated vertices), which glue back
along the cut vertices into the \term{block-cut tree}.  A \term{functional graph}
-- a set \(f\colon V\to V\), every vertex one outgoing arrow -- has each weak
component shaped like a \term{rho}: one directed cycle with ordinary rooted trees
hanging into it.  A \term{cactus} is a graph where every edge lies in at most one
simple cycle, so its cycles are exactly its blocks and one contracted node per
cycle gives back a tree.  In each case the graph is not a tree, but a tree is
one DFS away, and \emph{that} tree is what you count, colour, DP and answer
queries on -- the content of Chapters~12 (\chref{c09}), 13 (\chref{c10}) and 17
(\chref{c28}) becomes legal the moment you have built it.

\textbf{You will learn to}
\begin{itemize}[nosep]
  \item find every bridge in \(\bigO{n+m}\) with the low-link test
        \(\mathrm{low}[v]>\tin[u]\), contract the 2-edge-connected components, and
        read the bridge tree (\(\secref{c31-bridges}\), solved CF 118E \emph{Bertown
        roads}: orient all edges for strong connectivity iff no bridge exists); \hfill\lv{4}
  \item find every articulation point and pop each \emph{block} off an edge stack
        when \(\mathrm{low}[v]\ge\tin[u]\), then answer edge-level cycle questions
        on the block-cut tree (\(\secref{c31-block}\), solved CF 962F \emph{Simple
        Cycles Edges}: an edge is on exactly one simple cycle iff its block has
        \(|E|=|V|\)); \hfill\lv{5}
  \item decompose a functional graph into its rho pieces -- peel in-degrees to
        find the cycle, then run Ch.13 DP on each in-tree (\(\secref{c31-functional}\)); \hfill\lv{4}
  \item contract a cactus along its cycles and DP over the resulting tree with a
        cycle-transfer at each cycle node (\(\secref{c31-cactus}\)); \hfill\lv{4}--\lv{5}
  \item read a graph statement and pick bridge-tree vs.\ block-cut vs.\ functional
        vs.\ cactus -- and know when the graph has no useful tree at all
        (\(\secref{c31-boundary}\)). \hfill\lv{5}
\end{itemize}

\textbf{Prerequisites.} Chapter~12 (\chref{c09}) -- diameter and height by two
sweeps: the black box the Summit runs on each attached tree.  Chapter~13
(\chref{c10}) -- the bottom-up fold: what you run ON the extracted tree.
Chapter~17 (\chref{c28}) -- DSU / component merging: the contraction engine
behind the bridge tree.  Chapter~5 (\chref{c24}) -- \(\tin/\tout\): the same
DFS-order the low-link test reads.  Chapter~28 (\chref{c22}) -- the
\textbf{previous chapter}; it ends where this one begins (``the input is a
tree'' \(\to\) ``here it is not'').
\end{chaptertheorybox}

% ---------------- Reading paths + ladder position ----------------
\begin{chaptertheorybox}{How to Read This Chapter (three paths)}
Level badges \(\lv3\)--\lv{6}: \(\lv3\) is running the low-link bridge test by
hand on the running graph and naming the bridges; \(\lv4\) is coding the
contraction and the CF 118E orientation, and reading a functional graph as rho
pieces; \(\lv5\) is the edge-stack block pop (962F) and the \(|E|=|V|\) cycle
lemma; \(\lv6\) is the Summit (CF 835F) -- a unicyclic graph, one cycle-edge
deletion, a sliding-window diameter.

\textbf{Edge-connectivity path (\(\lv3\)--\lv{4}, ``cut edges'').}
\secref{c31-bridges} in full -- tin/low, the bridge test, contract to the bridge
tree -- capped by solved CF 118E, where Robbins' theorem turns ``orient for
strong connectivity'' into ``check for a bridge''.  Stop here if you only need
edge-level answers.

\textbf{Vertex-connectivity path (\(\lv4\)--\lv{5}, ``cut vertices'').}
\secref{c31-block} -- articulation points, the edge-stack block pop, the
block-cut tree -- then solved CF 962F, where the whole question is which
\emph{blocks} are simple cycles.  The \(|E|=|V|\) lemma is brute-force verified;
read the ledger before trusting it.

\textbf{Graphs-that-are-secretly-trees path (\(\lv4\)--\lv{6}).}
\secref{c31-functional} (rho decomposition, in-degree peeling),
\secref{c31-cactus} (cycles-as-blocks, contract and DP), then
\(\secref{c31-boundary}\) for the honest limits -- and the Summit
\(\probref{c31-summit}\), a unicyclic graph (a cactus with one cycle) where the
tree you extract is a \emph{cycle of trees}, and the answer is a diameter over a
sliding window on it.
\end{chaptertheorybox}

\begin{figure}[H]
\centering
\begin{tikzpicture}[diagram, x=8.5mm, y=6.5mm]
  \foreach \st/\lab in {1/lv1,2/lv2,3/lv3,4/lv4,5/lv5,6/lv6} {
    \fill[\lab!22, draw=\lab, line width=0.8pt] (\st-1,0) rectangle (\st,\st*0.52);
    \node[font=\scriptsize\bfseries, text=\lab] at (\st-0.5,\st*0.52-0.42) {L\st};}
  \node[annot, anchor=south west, font=\scriptsize]
    at (0.05,0.62) {prereq (Ch.12/13/17): diameter, bottom-up fold, DSU};
  \node[annot, anchor=south west, font=\scriptsize]
    at (2.05,1.66) {29.1 tin/low bridges, contract the bridge tree -- CF 118E};
  \node[annot, anchor=south west, font=\scriptsize]
    at (3.05,2.18) {29.3 functional rho | 29.4 cactus contraction};
  \node[annot, anchor=south west, font=\scriptsize, text=lv5]
    at (4.05,2.70) {29.2 edge-stack blocks -- CF 962F (|E|=|V| cycle blocks)};
  \node[annot, anchor=south west, font=\scriptsize, text=lv6]
    at (5.05,3.22) {29.5 the pick; Summit CF 835F (unicyclic min-diameter)};
\end{tikzpicture}
\caption{Chapter~29 on the six-step ladder: hand-running the bridge test on the
running graph is \(\lv3\); the coded contraction and CF 118E, and reading
functional/cactus graphs, are \(\lv4\); the edge-stack block pop and CF 962F are
\(\lv5\); the Summit (CF 835F, delete one cycle edge to minimise the tree
diameter) is step~6.}
\end{figure}

% ---------------- The running example ----------------
\subsection*{The Running Example: A Graph With Two Cycles and a Tail}

Every section runs on the same 9-vertex undirected graph \(G\), fixed once here
and authoritative (if a later trace disagrees, this figure wins).  Edges, 1-based:
\[
  1\!-\!2,\; 2\!-\!3,\; 3\!-\!1,\quad
  3\!-\!4,\quad
  4\!-\!5,\; 5\!-\!6,\; 6\!-\!4,\quad
  4\!-\!7,\; 7\!-\!8,\; 8\!-\!9 .
\]
Two triangles meet at a single bridge: \(B_1=\{1,2,3\}\) (edges
\(1\!-\!2,2\!-\!3,3\!-\!1\)) and \(B_2=\{4,5,6\}\) (edges
\(4\!-\!5,5\!-\!6,6\!-\!4\)), joined by the bridge \(3\!-\!4\); from vertex \(4\)
a pendant path \(4\!-\!7\!-\!8\!-\!9\) hangs, every one of its three edges a
bridge.  The \textbf{bridges} are exactly \(\{3\!-\!4,\;4\!-\!7,\;7\!-\!8,\;8\!-\!9\}\);
the \textbf{articulation points} are \(\{3,4,7,8\}\).  Delete all bridges and the
\term{2-edge-connected components} are \(B_1,B_2,\{7\},\{8\},\{9\}\), so the
\term{bridge tree} is the path \(B_1\!-\!B_2\!-\!7\!-\!8\!-\!9\).  Cut along the
articulation points instead and the \term{blocks} are the two triangles plus the
four bridges; the \term{block-cut tree} interleaves block-nodes and cut-vertices.

\begin{figure}[H]
\centering
\begin{tikzpicture}[diagram, scale=0.95, transform shape,
   tnode/.style={circle, draw=primaryblue!70, fill=primaryblue!8, line width=0.7pt,
                 minimum size=6.6mm, font=\scriptsize\bfseries, inner sep=1pt},
   cut/.style={circle, draw=difforange, fill=difforange!18, line width=1.1pt,
             minimum size=6.6mm, font=\scriptsize\bfseries, inner sep=1pt}]
  \node[tnode] (v1) at (0,1.7) {1};
  \node[tnode] (v2) at (-1.2,0.2) {2};
  \node[cut]   (v3) at (1.2,0.2) {3};
  \node[cut]   (v4) at (3.2,0.2) {4};
  \node[tnode] (v5) at (2.3,-1.4) {5};
  \node[tnode] (v6) at (4.1,-1.4) {6};
  \node[cut]   (v7) at (5.3,0.2) {7};
  \node[cut]   (v8) at (6.6,0.2) {8};
  \node[tnode] (v9) at (7.6,-0.9) {9};
  \draw[primaryblue!70, line width=1.1pt] (v1)--(v2) (v2)--(v3) (v3)--(v1);
  \draw[primaryblue!70, line width=1.1pt] (v4)--(v5) (v5)--(v6) (v6)--(v4);
  \draw[warningred, line width=1.3pt] (v3)--(v4) (v4)--(v7) (v7)--(v8) (v8)--(v9);
  \node[annot, font=\scriptsize, anchor=north west, text width=8.5cm] at (-1.6,2.4)
    {blue = inside a 2-edge-connected block;  red = a bridge;  orange = an articulation point};
\end{tikzpicture}
\caption{The running graph \(G\): triangles \(B_1=\{1,2,3\}\) and \(B_2=\{4,5,6\}\)
in blue, joined by the red bridge \(3\!-\!4\), with the red pendant path
\(4\!-\!7\!-\!8\!-\!9\).  Articulation points (orange) are \(3,4,7,8\).  The bridge
tree contracts \(B_1,B_2,\{7\},\{8\},\{9\}\) to the path \(B_1\!-\!B_2\!-\!7\!-\!8\!-\!9\);
the block-cut tree is traced in \(\secref{c31-block}\).}
\label{fig:c31-graph}
\end{figure}

\begin{figure}[H]
\centering
\begin{tikzpicture}[diagram, scale=1.0, transform shape,
   comp/.style={rectangle, rounded corners=2pt, draw=primaryblue!70,
                fill=primaryblue!8, line width=0.8pt, minimum size=8mm,
                font=\scriptsize\bfseries, inner sep=2pt}]
  \node[comp] (b1) at (0,0) {\(B_1\)};
  \node[comp] (b2) at (2,0) {\(B_2\)};
  \node[comp] (c7) at (3.6,0) {\(\{7\}\)};
  \node[comp] (c8) at (5,0) {\(\{8\}\)};
  \node[comp] (c9) at (6.4,0) {\(\{9\}\)};
  \draw[warningred, line width=1.3pt] (b1)--(b2) (b2)--(c7) (c7)--(c8) (c8)--(c9);
  \node[annot, font=\scriptsize] at (3.2,-0.85)
    {each red edge is one bridge of \(G\)};
\end{tikzpicture}
\caption{The bridge tree of \(G\): contract every 2-edge-connected component to a
node; the bridges become the tree edges.  It is a tree precisely \emph{because}
every remaining connection was a bridge -- the object Chapters 12/13 let you DP on.}
\label{fig:c31-bridgetree}
\end{figure}

Two small traps return throughout.  The \textbf{bridge-is-not-cut-vertex trap}:
an articulation \emph{point} and a bridge \emph{edge} are different failures.  In
\(G\) removing vertex \(3\) disconnects \(\{1,2\}\) from the rest, yet no single
edge incident to \(3\) inside the triangle \(B_1\) is a bridge -- the triangle
survives any one edge.  Conversely an endpoint of a bridge need not be a cut
vertex: \(9\) is incident to the bridge \(8\!-\!9\) but removing \(9\) disconnects
nothing.  Edge-connectivity and vertex-connectivity are separate ledgers
(\(\secref{c31-bridges}\) vs.\ \(\secref{c31-block}\)).  The
\textbf{block-is-not-a-2-edge-component trap} is subtler: \(B_1\) is both a
2-edge-connected component \emph{and} a block, but the pendant edges
\(4\!-\!7,7\!-\!8,8\!-\!9\) are each their own block (a bridge is a block of one
edge) while \(\{7\},\{8\}\) are trivial 2-edge-components -- so the block-cut tree
has four bridge-blocks and a different shape from the bridge tree.  Confusing the
two decompositions is the single most common wrong answer on CF 962F.

\subsection*{Which Job? Reading the Statement}

\begin{center}
\renewcommand{\arraystretch}{1.25}
\setlength{\tabcolsep}{3pt}
\begin{tabular}{@{}p{0.30\linewidth} p{0.20\linewidth} p{0.28\linewidth} >{\raggedright\arraybackslash}p{0.13\linewidth}@{}}
\toprule
\textbf{The statement says\dots} & \textbf{Extract} & \textbf{Why it is a tree} & \textbf{Section}\\
\midrule
``removing this edge disconnects'' / ``orient every edge'' & bridge tree & bridges are exactly the tree edges of the contraction & \(\secref{c31-bridges}\)\\
``removing this vertex disconnects'' / ``on exactly one cycle'' & block-cut tree & blocks meet only at cut vertices \(\to\) a bipartite tree & \(\secref{c31-block}\)\\
``each \(x\) points to \(f(x)\)'' / ``eventually cycles'' & functional rho & out-degree 1 \(\Rightarrow\) one cycle + in-trees per component & \(\secref{c31-functional}\)\\
``each edge in at most one cycle'' & cactus contraction & cycles-as-blocks, contract each to a node & \(\secref{c31-cactus}\)\\
\bottomrule
\end{tabular}
\end{center}

\begin{insightbox}[The whole chapter in one sentence]
A general graph is a tree \emph{plus} a set of redundant edges; kill the redundancy
(find the bridges, the cut vertices, the functional cycles) and what is left is a
tree you already know how to work on -- the entire craft of this chapter is
\emph{choosing which redundancy to kill} so that the surviving tree answers the
question actually asked.
\end{insightbox}

% ---------------- Taxonomy figure ----------------
\begin{figure}[H]
\centering
\begin{tikzpicture}[diagram, scale=0.9, transform shape,
   mnode/.style={circle, draw=primaryblue!70, fill=primaryblue!8, line width=0.7pt,
                 minimum size=13mm, font=\scriptsize\bfseries, inner sep=1pt, align=center}]
  \node[mnode] (root) at (0,0) {graph};
  \node[mnode] (dir)  at (-4,-1.6) {out-deg 1\\ function};
  \node[mnode] (un)   at (1.5,-1.6) {undirected};
  \node[mnode] (fg)   at (-5.6,-3.2) {\(\rho\) pieces};
  \node[mnode] (fr)   at (-2.6,-3.2) {bridge\\ tree};
  \node[mnode] (bc)   at (0.4,-3.2) {block-cut\\ tree};
  \node[mnode] (cac)  at (3.4,-3.2) {cactus};
  \draw (root)--(dir)--(fg);
  \draw (root)--(un)--(fr);
  \draw (un)--(bc)--(cac);
\end{tikzpicture}
\caption{Where the four decompositions sit. A functional graph is already
directed and out-degree one \(\to\) rho pieces (\(\secref{c31-functional}\)).
An undirected graph: kill \emph{edge} redundancy \(\to\) bridge tree
(\(\secref{c31-bridges}\)); kill \emph{vertex} redundancy \(\to\) block-cut tree
(\(\secref{c31-block}\)); if cycles never share an edge the two coincide into the
cactus (\(\secref{c31-cactus}\)).}
\label{fig:c31-taxonomy}
\end{figure}

\subsection*{Who Owns What}

\begin{center}
\renewcommand{\arraystretch}{1.2}
\setlength{\tabcolsep}{3pt}
\begin{tabular}{@{}p{0.28\linewidth} p{0.40\linewidth} p{0.24\linewidth}@{}}
\toprule
\textbf{Object} & \textbf{Owned here?} & \textbf{Elsewhere}\\
\midrule
Bridge / 2-edge-component & \textbf{Yes}, \(\secref{c31-bridges}\) -- find and contract & --\\
Block-cut tree, articulation & \textbf{Yes}, \(\secref{c31-block}\) & --\\
Diameter / height of a piece & no -- used as a black box & Ch.12 (\chref{c09})\\
Bottom-up DP on the tree & no -- the fold is reused & Ch.13 (\chref{c10})\\
Component merge (contraction) & no -- DSU is a black box & Ch.17 (\chref{c28})\\
Minimum spanning tree & no -- a different ``tree in a graph'' & Ch.31 (\chref{c17})\\
\bottomrule
\end{tabular}
\end{center}

\subsection*{Recognition Guide}

\begin{recognitionbox}[Turning a graph into a tree you can DP on]
\reckey{``disconnect'', ``removing one edge/vertex'', ``still connected after any
\(k\) removals'', ``on exactly one cycle'', ``every \(x\) maps to \(f(x)\)``,
``each edge in at most one cycle'', ``is it a tree / cactus?''}
\recuse{
  \item \textbf{Answer depends only on which edge-cuts are unavoidable} -- contract
        2-edge-connected components and DP on the bridge tree (\(\secref{c31-bridges}\)).
  \item \textbf{Answer depends on vertices shared by cycles} -- pop blocks off the
        edge stack and use the block-cut tree (\(\secref{c31-block}\)).
  \item \textbf{Out-degree exactly one} -- peel in-degrees, find the rho, DP on the
        in-trees (\(\secref{c31-functional}\)).
}
\recnotuse{
  \item \textbf{You need a minimum total weight} -- that is a spanning tree, not a
        decomposition: Ch.31 (\chref{c17}).
  \item \textbf{Redundancy is directional reachability, not edge/vertex cuts} -- that
        is strong-connectivity / dominators, largely graph-land; \(\secref{c31-boundary}\)
        gives the honest boundary.
}
One-line test: \emph{if deleting a specified edge or vertex can change the answer,
a bridge/articulation decomposition is on the table; build the tree, then reuse
Ch.12/13.}
\end{recognitionbox}

\subsection*{Constraint \(\to\) Technique}

\begin{center}
\renewcommand{\arraystretch}{1.2}
\setlength{\tabcolsep}{3pt}
\begin{tabular}{@{}p{0.34\linewidth} p{0.34\linewidth} >{\raggedright\arraybackslash}p{0.20\linewidth}@{}}
\toprule
\textbf{Constraint signature} & \textbf{Move} & \textbf{Section}\\
\midrule
\(n,m\le10^5\), ``orient / 2-edge-connected'' & low-link bridges \(\to\) contract & \(\secref{c31-bridges}\)\\
``exactly one cycle'' per edge / block queries & edge-stack blocks \(\to\) \(|E|=|V|\) & \(\secref{c31-block}\)\\
given a function \(f\), ``who reaches a cycle'' & in-degree peel \(\to\) rho DP & \(\secref{c31-functional}\)\\
connected, \(n\) edges (unicyclic), minimise diameter & find cycle, window the trees & \(\probref{c31-summit}\)\\
\bottomrule
\end{tabular}
\end{center}

\subsection*{Notation Used Only Here}

\(\mathrm{tin}[v]\) is the DFS entry time (Ch.5, \chref{c24});
\(\mathrm{low}[v]\) is the smallest \(\mathrm{tin}\) reachable from the subtree of
\(v\) using at most one back edge.  An edge \(u\!-\!v\) (tree edge, \(v\) the
child) is a \term{bridge} iff \(\mathrm{low}[v]>\mathrm{tin}[u]\).  A
\emph{block} is a maximal set of edges any two of which lie on a common simple
cycle (equivalently, a maximal 2-vertex-connected subgraph, with bridges and
isolated vertices as degenerate blocks); a \emph{cut vertex} is a vertex in two or
more blocks.  For a functional graph, \(\rho(v)\) is the length of the tail from
\(v\) to the cycle it feeds and \(\mu(v)\) the cycle length.  The Summit reads
\(h[c]\) and \(\mathrm{dtree}[c]\) for the height and internal diameter of the
in-tree hanging off cycle vertex \(c\).
